\documentclass{article}

\usepackage{amsmath,amssymb}
\usepackage{xurl}
\usepackage[hidelinks]{hyperref}
\setlength{\emergencystretch}{2em}

% Shared commands for notes in docs/paper-gaps/.

\DeclareUnicodeCharacter{2080}{\ifmmode{}_{0}\else\textsubscript{0}\fi}
\DeclareUnicodeCharacter{2081}{\ifmmode{}_{1}\else\textsubscript{1}\fi}
\DeclareUnicodeCharacter{2082}{\ifmmode{}_{2}\else\textsubscript{2}\fi}
\DeclareUnicodeCharacter{2099}{\ifmmode{}_{n}\else\textsubscript{n}\fi}

\newcommand{\Lean}{\textsc{Lean}}
\ExplSyntaxOn
\NewDocumentCommand{\leanid}{m}{
  \ifmmode
    \textsf{\detokenize{#1}}
  \else
    \group_begin:
    \urlstyle{sf}
    \tl_set:Nn \l_tmpa_tl {#1}
    \tl_replace_all:Nnn \l_tmpa_tl {\_} {_}
    \exp_args:NV \path \l_tmpa_tl
    \group_end:
  \fi
}
\ExplSyntaxOff
\providecommand{\tr}{\operatorname{tr}}
\newcommand{\tnleanIssueBase}{https://github.com/LionSR/TNLean/issues/}
\newcommand{\tnleanPrBase}{https://github.com/LionSR/TNLean/pull/}
\newcommand{\tnleanDocBase}{https://sirui-lu.com/TNLean/docs/}
\newcommand{\ghissue}[1]{\href{\tnleanIssueBase#1}{GitHub issue \##1}}
\newcommand{\ghpr}[1]{\href{\tnleanPrBase#1}{PR \##1}}
\newcommand{\doclink}[1]{\href{\tnleanDocBase#1.html}{\path{#1}}}

% Dirac notation and the identity operator, as in blueprint/src/macros/common.tex.
% \providecommand keeps note-local definitions authoritative.
\usepackage{dsfont}
\providecommand{\ket}[1]{|#1\rangle}
\providecommand{\bra}[1]{\langle #1|}
\providecommand{\braket}[2]{\langle #1 | #2 \rangle}
\providecommand{\ketbra}[2]{|#1\rangle\!\langle #2|}
\providecommand{\Id}{\mathds{1}}
\providecommand{\one}{\mathds{1}}
\providecommand{\C}{\mathbb{C}}
\providecommand{\MN}[1]{M_{#1}(\C)}
\providecommand{\MD}{M_D(\C)}

% Verdict marker: \gapnote{<kind>}{<status>}. Kind is the policy
% classification (clarification, local-correction, scope-restriction,
% unfaithful, false-source, open-gap); status is open, wip (elimination
% underway), resolved, or historical. Machine-read by the site index;
% prints a verdict line.
\newcommand{\gapnote}[2]{\par\noindent\textbf{Verdict:} #1 (#2).\par}


% Fallbacks keep this note compilable when a different command.tex is found first.
\providecommand{\leanid}[1]{\nolinkurl{#1}}
\providecommand{\gapnote}[2]{\par\noindent\textbf{Verdict:} #1 (#2).\par}

\title{The Finite-Range Approximation of a Filtered Term in Hastings' Local
  Solution of Gapped Hamiltonians}
\author{}
\date{2026-10-23}

\begin{document}
\maketitle
\gapnote{local-correction}{resolved}

\section{The source passage}

In~\cite{Hastings2006Solving} (arXiv:cond-mat/0508554v3, \texttt{proj3new.tex})
the Hamiltonian is $H=\sum_iH_i$ on a lattice with a metric, with
$\lVert H_i\rVert\le J$ and $H_i$ acting within distance $R$ of the site $i$
(lines~169--190). For $q>0$ and a gap $\Delta E>0$, the filtered term is
\begin{align}
  \tilde H_i^0=\frac{\Delta E}{\sqrt{2\pi q}}\int dt\,
  e^{iHt}H_ie^{-iHt}\,e^{-(t\Delta E)^2/(2q)}
  \label{eq:hastings05_filtered_term}
\end{align}
(eqs.~\texttt{tnd}, \texttt{thd}, lines~194--207). The source then defines
(eqs.~\texttt{mid}, \texttt{odef}, \texttt{hpdef}, lines~285--307)
\begin{align}
  M_i=\frac{\Delta E}{\sqrt{2\pi q}}\int dt\,
  e^{-(t\Delta E)^2/2}\,e^{iH_{\mathrm{loc}}t}H_ie^{-iH_{\mathrm{loc}}t},
  \qquad
  H_{\mathrm{loc}}=\sum_{j:\,d(i,j)\le l_{\mathrm{proj}}-R}H_j,
  \label{eq:hastings05_local_evolution_term}
\end{align}
and states that $M_i$ acts within distance $l_{\mathrm{proj}}$ of $i$, that
$\lVert M_i\rVert\le J$ (line~320), and that
\begin{align}
  \lVert\tilde H_i^0-M_i\rVert\le
  J\bigl(N(l_{\mathrm{proj}})\,g(c_1l_{\mathrm{proj}},l_{\mathrm{proj}})
  +2e^{-(c_1l_{\mathrm{proj}}\Delta E)^2/(2q)}\bigr)
  \label{eq:hastings05_local_evolution_bound}
\end{align}
(eq.~\texttt{apb}, lines~309--320, proved at lines~1458--1533). The later
argument (eq.~\texttt{facld} and lines~359--387) uses only these three
properties of $M_i$.

\section{Two local corrections}

\paragraph{The exponent in~\eqref{eq:hastings05_local_evolution_term}.}
The Gaussian weight in~\eqref{eq:hastings05_local_evolution_term} is printed as $e^{-(t\Delta E)^2/2}$. With the prefactor
$\Delta E/\sqrt{2\pi q}$ it is a probability density only when $q=1$, and the
comparison~\eqref{eq:hastings05_local_evolution_bound}
with~\eqref{eq:hastings05_filtered_term} splits the two integrals at the same
time with the same weight. The weight of~\eqref{eq:hastings05_filtered_term},
$e^{-(t\Delta E)^2/(2q)}$, is meant.

\paragraph{The construction of $M_i$.} The formalization does not evolve $H_i$
under $H_{\mathrm{loc}}$. It takes for $M_i$ the normalized partial trace
$\mathbb E_K(\tilde H_i^0)$ of the filtered term onto the set $K$ of sites
within graph distance $l$ of the support of $H_i$: the trace over the sites
outside $K$, divided by their dimension, tensored with the identity there.
This operator has analogous support, norm and approximation bounds.
\begin{itemize}
  \item It acts on $K$.
  \item $\lVert\mathbb E_K(\tilde H_i^0)\rVert\le\lVert\tilde H_i^0\rVert
    \le\lVert H_i\rVert\le J$, since the partial trace and the filter do not
    increase the norm. It is Hermitian when $H_i$ is.
  \item Splitting~\eqref{eq:hastings05_filtered_term} at $|t|=T$ and applying
    the Lieb--Robinson localization of $e^{iHt}H_ie^{-iHt}$ onto $K$ for
    $|t|\le T$ gives
    \begin{align}
      \lVert\tilde H_i^0-\mathbb E_K(\tilde H_i^0)\rVert
      \le C'|S|^2\lVert H_i\rVert e^{vT-l/2}
      +2\lVert H_i\rVert e^{-(T\Delta E)^2/(2q)},
      \label{eq:hastings05_filtered_partial_trace_bound}
    \end{align}
    with $S$ the support of $H_i$ and $C'$, $v$ the constants of the graph
    localization estimate. This is an analogue
    of~\eqref{eq:hastings05_local_evolution_bound}, with a different construction and localization term.
\end{itemize}

Since $S$ lies within distance $R$ of $i$, the set $K$ lies within distance
$l+R$ of $i$. For integer radii $l_{\mathrm{proj}}\ge R$, take
$l=l_{\mathrm{proj}}-R$ to recover the source's support radius. The
localization term in~\eqref{eq:hastings05_filtered_partial_trace_bound}
then becomes
\begin{align}
  C'|S|^2\lVert H_i\rVert e^{vT-(l_{\mathrm{proj}}-R)/2}.
  \label{eq:hastings05_filtered_shifted_localization}
\end{align}
In particular, at $T=c_1l_{\mathrm{proj}}$ the
term~\eqref{eq:hastings05_filtered_shifted_localization} carries an extra factor
$e^{R/2}$ compared with substituting $l=l_{\mathrm{proj}}$. The source's
quantitative estimate is not asserted for this different construction.
At $q=0$ the filter is defined to be $H_i$ itself; the bounds remain valid
with the Gaussian tail factor replaced by $1$, without division by zero.

\section{Formalization}

\leanid{TNLean.Lattice.exists_local_approx_gaussianFilteredTerm} gives an
operator with these support, norm and approximation bounds, and
\leanid{TNLean.Lattice.norm_gaussianFilteredTerm_sub_siteExpectation_le_graphNeighborhood}
proves~\eqref{eq:hastings05_filtered_partial_trace_bound} for the partial
trace. The theorems that take $M_i$ as input,
\leanid{TNLean.Lattice.norm_sub_starProjection_apply_le_of_gaussianFilteredTerm},
\leanid{TNLean.Lattice.norm_apply_sub_inner_smul_le_of_gaussianFilteredTerm}
and \leanid{TNLean.Lattice.norm_sum_sub_le_of_sum_eq}, hold for every operator
$M$, so they apply to the source's truncated evolution as well.

\bibliographystyle{plain}
\bibliography{references}

\end{document}
